\paragraph{Inertia and the abelian quotients of the fitted frame.}
The generator count uses these inertia quotients.

Let \(G\) be the quotient of \(U_m=G_{L^+}\) by the closed
normal subgroup generated by the away-from-\(\ell\) inertia
elements on which both \(\eta\) and \(P\) are one. Its positive
component is the image of \(G_L\). On the away-from-\(\ell\)
inertia of \(U_m\), the full character \(\Psi\) has prime-to-
\(\ell\) order, since \(H_b\) kills the other part and the
cyclotomic character is unramified there. Reduction on that
prime-to-\(\ell\) subgroup is injective. Thus the geometric
inertia condition on \(G\) kills the same kernel as the
ordered residual character at these places.

For each place of \(L^+\) above \(\ell\), use a representative
conjugate to \(B\) by an element of \(\ker\eta\). Let \(g\) be any conjugating element. If its sign is negative, multiply on the left by complex conjugation \(c\in U_m\):
\[
 \eta(g)=-1\quad\Longrightarrow\quad \eta(cg)=(-1)(-1)=1 .
\]
This preserves the \(U_m\)-orbit of the representative. The
places split in \(L/L^+\), so the chosen decomposition group
of \(U_m\) is also the decomposition group on the indicated
side over \(L\). These are the local groups mapped into the
frame. They are the actual decomposition groups before the
projection to \(G\).

There is a nontrivial mod-\(\ell\) cyclotomic value in each
of these local groups. On rational inertia at \(\ell\), the
quadratic character \(\eta\) is trivial in both host cases,
and so is the character of \(H_0\) by the local comparison.
Start with a nontrivial mod-\(\ell\) cyclotomic value and take
an \(\ell\)-power to land in \(H_b\cap H_w\cap H_m\).
Since the nontrivial residue value \(a\) lies in \(\mathbb F_\ell^\times\), raising it to an \(\ell\)-power preserves its exact order:
\[
 \operatorname{ord}(a^{\ell^t})=\operatorname{ord}(a)>1 .
\] Conjugating supplies
the element at every selected place.

In particular the
completions on that side contain no primitive \(\ell\)-th
root of unity. There is also a global element of \(G_L\)
on which \(P=1\) but \(\bar\epsilon\ne1\). Take the element
\(s\) of the preceding separation test, with
\((\eta,P,\bar\epsilon)(s)=(-1,1,1)\). At the product of \(s\) with complex conjugation the three character values are \((1,1,-1)\).
If needed, correct the character of \(H_0\) by an inertia
element at its new ramified prime \(p\). All three characters are trivial on this inertia group. An \(\ell\)-power again
lands in the remaining three kernels. This proves the
global separation required at the auxiliary level.

\begin{lemma}[The fitted unramified character bounds]
\label{lem:fitted-abelian-quotient}
In the frame of Proposition~\ref{prop:separated-marked-host}, let a
finite abelian \(\ell\)-group character of the positive component
vanish on \(P=1\) inertia at \(\ell\) and all its conjugates, including
complex conjugates. Its restriction \(\chi|_{G_L}\) factors as 
\[
 G_L\longrightarrow\operatorname{Cl}(L)[\ell^\infty]\longrightarrow\operatorname{im}\chi.
\]
 On the Gaussian sign kernel
with vanishing on the marked inertia intersections the unramified
pullback is trivial. For residual coefficients formed from
\(\eta,P,\bar\epsilon\), a class whose restrictions kill inertia
over \(E_m=E(\mu_{\ell^{m+1}})\) restricts to a character of
\(\operatorname{Cl}(E_m)\).
\end{lemma}

\begin{proof}
\proofstep{The unramified quotient and the class number.}
Consider a continuous finite abelian \(\ell\)-group character on
the positive component of \(G\). Suppose it vanishes on the
indicated \(P=1\) inertia at \(\ell\) and its conjugates,
including the conjugates by complex conjugation. The quotient
of inertia by that kernel has order dividing \(\#k-1\), where
\(k\) is the residue field. This quotient has trivial image
in the \(\ell\)-group. At every place away from \(\ell\),
the inertia of \(G_L\) already has image of order dividing
\(\#k-1\), by the definition of \(G\). The resulting character on \(G_L\) is unramified with \(\ell\)-primary image.

By class field theory its factorization at this stage goes through
\[
 G_L\longrightarrow\operatorname{Cl}(L)[\ell^\infty]\longrightarrow\operatorname{im}(\chi),
\] Its exponent gives one
positive annihilating integer for every truncation of a
principal-unit character. An involution-invariant character
with the specified vanishing on inertia on one side has
it on the other side as well. This is the case for the
middle character in the exceptional corner. On the original
Gaussian quotient \(G_{\Q,S}\), use instead the sign kernel
giving \(G_{\Q(i)}\). A character trivial on the conjugates
of all the indicated intersections with inertia pulls back
to an everywhere unramified character of \(\Q(i)\). The norm-Euclidean property of the Gaussian integers makes every ideal principal. Hence
\[
 \operatorname{Cl}(\mathbb Q(i))=1,\qquad h_{\mathbb Q(i)}=1 .
\]
Here the pullback character is trivial at each finite
truncation. \cite{ray-class}\label{cor:ordinary-fitted-arithmetic-1}

\proofstep{Restriction of residual classes.}
The field \(E_m=E(\mu_{\ell^{m+1}})\) is
obtained from \(L^+\) by adjoining the fields of
\(\eta,P,\bar\epsilon\), of degrees prime to \(\ell\).
For any residue character formed from these three, restriction
of \(H^1\) to \(G_{E_m}\) is injective by
inflation--restriction and prime-to-\(\ell\) averaging.
There the coefficients are constant. In particular when the
local conditions kill inertia on this subgroup, the classes
embed in the additive character group of the class group of
\(E_m\). 
\end{proof}

\paragraph{Uniform generator bounds and local deformation data.}
With the abelian quotients fixed, we can count generators and record the local
conditions for the relative cohomology calculation.

By Lemma~\ref{lem:fitted-abelian-quotient}, the strict residual
classes over \(E_m\) embed in the additive character group of its
class group. Its dimension over \(k\) is bounded by \(C_0\).
Apply this to the strict group for \(\eta P^{-2}\) over
\(L^+\) and to that for \(P^{-1}\) over \(L\). In the
first case impose a coboundary at every marked local
group.

In the second impose it at both conjugate places.
The local character is nontrivial in each case by
\(P(g_*)^2\ne1\), transported to the marked local groups.
Indeed, if \(b,b'\) trivialize the same local coboundary for the nontrivial character \(\chi\), then for an element \(g\) with \(\chi(g)\ne1\),
\[
 (\chi(g)-1)(b-b')=0\quad\Longrightarrow\quad b=b' .
\]
Keeping these scalars in the relative complex adds no
dimension to its first cohomology. After restriction to
\(G_{E_m}\), a coboundary is zero. At the other finite
places inertia below \(G_{E_m}\) is killed by the frame
quotient. The same two bounds by \(C_0\) therefore hold for
the relative strict groups.

\begin{lemma}[Additive character bounds and deformation directions]
\label{lem:fitted-local-generators}
Use the fitted frame quotient of Lemma~\ref{lem:fitted-abelian-quotient}.
For \(\eta P^{-2}\) over \(L^+\) require a coboundary at every marked
local group, and for \(P^{-1}\) over \(L\) require it at both conjugate
places. Impose unramified conditions at the remaining finite places.
Writing \(S_{L^+}(\eta P^{-2})\) and \(S_L(P^{-1})\) for the resulting
relative strict character spaces, one has
\[
  \max\bigl(\dim S_{L^+}(\eta P^{-2}),\ \dim S_L(P^{-1})\bigr)\le C_0.
\]
For the trivial character use the full local conditions at \(\ell\)
and the unramified conditions elsewhere. The ordinary local deformation
directions can be chosen independently at the marked places.
\end{lemma}

\begin{proof}
\proofstep{The bound for additive characters.}
For the trivial character, every additive character on
\(G\) is unramified away from \(\ell\). Take the Selmer group
over \(L^+\) with full conditions at \(\ell\) and unramified
conditions elsewhere. The dual coefficient line is
\(k(1)\). Its global invariant space is zero by the
cyclotomic test. Its local invariant spaces at \(\ell\) are
also zero. The Greenberg--Wiles formula gives the difference
of the two Selmer dimensions as
\[
           1+\sum_{v\mid\ell} d_v-[L^+:\Q]=1 .
\]
Indeed the local full condition contributes \(d_v\) by
Tate duality and the local Euler characteristic. An
unramified finite-place condition contributes zero.

At each real place the invariant dimension in the trivial
line is one. The dual Selmer group has zero local condition
at \(\ell\) and is unramified at the remaining finite
places, by orthogonality of unramified conditions away
from \(\ell\). Restriction to \(G_{E_m}\) bounds it by
\(C_0\). This gives the \(1+C_0\) bound for the first
cohomology with trivial coefficients. All the bounds use
the same frame as the earlier locally split cocycles.
The class rank bound precedes the choice of the deep
layer. If \(g\) is the number of these local places, summing the degree lower bounds gives
\[
 g\ell^m\le\sum_{v\mid\ell}[L^+_v:\mathbb Q_\ell]=[L^+:\mathbb Q]=d .
\]
\cite{ferrero-washington,washington-cyclo,milne-duality}

The Hecke levels use a second group, the Galois group
unramified outside the level set over \(L^+\), and its
quadratic kernel. They map to the fitted frame since the
closed inertia subgroup outside those levels is in the
kernel of the frame projection. Both sign and cyclotomic
character descend to the level group. For the determinant and patching constructions the finiteness condition on each level group and local group \(G\) is 
\[
 \dim_{\mathbb F_\ell}\operatorname{Hom}_{\rm cont}(U,\mathbb F_\ell)<\infty\qquad(U\subseteq G\text{ open}).
\]
 To see this for an open subgroup \(U\)
of a level group \(\Gamma\), apply global cohomological
finiteness to
\(\operatorname{Coind}_{U}^{\Gamma}\mathbb F_\ell\).
This is the finite continuous permutation module on
\(\Gamma/U\) (equivalently on the opposite cosets for the
right action). The level set is finite and includes the places
above \(\ell\). Shapiro gives
\[
 \operatorname{Hom}_{\rm cts}(U,\mathbb F_\ell)
   =H^1(U,\mathbb F_\ell)
   \simeq H^1(\Gamma,\operatorname{Coind}_{U}^{\Gamma}
                                    \mathbb F_\ell).
\]
The last group is finite by the global finiteness theorem.
This applies within the open quadratic kernel as well.

For a
local decomposition group use the same formula and local
finiteness for its completion. It also gives the finiteness
on the original rational level for the Gaussian base. If \(P_U\)
is the maximal pro-\(\ell\) quotient, these homomorphisms are the
dual of \(P_U/\overline{P_U^\ell[P_U,P_U]}\). The Frattini
criterion now makes \(P_U\) topologically finitely generated. This is the finite generation used in
the two universal determinant constructions.
\cite{milne-duality,chenevier}\label{cor:ordinary-fitted-arithmetic-2}

\proofstep{Independent deformation directions above the coefficient prime.}
At a selected \(\ell\)-place the maximal pro-\(\ell\) quotient
of the abelianized inertia image is, by local reciprocity,
the principal units of \(L_w\). At sufficiently large depth \(s\), logarithm identifies the deep units with an additive lattice of finite index in the integers:
\[
 \log(1+\mathfrak p_w^s)=\mathfrak p_w^s,\qquad \operatorname{rank}_{\mathbb Z_\ell}\mathfrak p_w^s=[L_w:\mathbb Q_\ell].
\]
The torsion is exactly the group of \(\ell\)-power roots of
unity, which is trivial by the cyclotomic test above.
Inside inertia the intersection of the kernels of \(P\) and
\(\bar\epsilon\) has prime-to-\(\ell\) index. Its image in
the free abelian pro-\(\ell\) unit group is the whole group,
since the cokernel has order dividing that index. On this
intersection the remaining cyclotomic map is the map
from the unit group to \(1+\ell\mathbb Z_\ell\). The image is
nonzero and open, by finite index in rational inertia.

If \(U_v\) denotes this pro-\(\ell\) unit lattice,
the remaining cyclotomic character gives an exact sequence
\[
 0\longrightarrow \ker\epsilon_\ell|_{U_v}
  \longrightarrow U_v\xrightarrow{\epsilon_\ell}
                   \epsilon_\ell(U_v)\longrightarrow1 .
\]
The image is free of rank one over \(\mathbb Z_\ell\),
so it splits by lifting one generator. The kernel is
therefore free of rank \(d_v-1\). Choose its basis along
with that one lift. Lift all \(d_v\) basis elements to the
intersection in inertia. Each has \(P=1\), and all but the
distinguished lift have full cyclotomic value one.
Vanishing on this basis makes a finite abelian
\(\ell\)-group character of the local group vanish on
inertia. This is the generator choice used to compare
two locally coincident labels on
\(N_v=I_v\cap\ker\epsilon\).
On the indicated inertia group \(P=\bar\epsilon^j\), so the defining kernel of \(N_v\) gives
\[
 P|_{N_v}=\bar\epsilon^j|_{N_v}=1 .
\]

For the decomposition group there is one further unramified
generator, obtained by lifting a uniformizer. For a
prescribed determinant weight at rank \(n\), divide each
inertia character by the Teichm\"uller lift of its residue.
Maps from the resulting free lattice to principal units
are specified by the values of a basis, which can be any
\(1+X_{vi}\) in the target. There are \(d_v\) free inertia coordinates for each of the \(n\) ordered characters at \(v\). With \(d=\sum_{v\mid\ell}d_v\), the weight algebra has the presentation 
\[
 \mathcal O[[X_1,\ldots,X_{nd}]],\qquad \sum_{v\mid\ell}n d_v=nd.
\]
 At a finite place away from \(\ell\) the
pro-\(\ell\) tame inertia quotient is procyclic. A chosen
lift of Frobenius together with inertia generates the
decomposition group topologically, since the unramified
quotient is generated densely by Frobenius. These statements
use inertia in the selected decomposition group, with
conjugacy transporting the generators to another choice of
place. \cite{ray-class,milne-duality}\label{cor:ordinary-fitted-arithmetic-3}

\end{proof}

\paragraph{The relative Euler inequality on the marked quotient.}
The generator bounds alone do not compare deformations with their obstructions. Here
the ordinary flags enter through the local quotients in a relative Euler
characteristic.

\begin{proposition}[Relative Euler inequality on the fitted quotient]
\label{prop:fitted-relative-euler}
Let \(\ell\) be odd. Let \(G\) be the quotient of \(G_{L^+}\) by
the closed normal subgroup generated by \(I_q\cap\ker\eta\cap\ker P\)
for \(q\nmid\ell\), and
let \(W\) be a finite-dimensional continuous representation of \(G\)
over a finite field \(k\) of characteristic \(\ell\). For each
\(v\mid\ell\) in \(L^+\), set \(d_v=[L_v^+:\mathbb Q_\ell]\)
and let \(b_v\subset W\) be stable under its marked group.
Assume \(\operatorname{Hom}_{G}(W,k(1))=0\).
Put \(Q_v=W/b_v\), and let \(H_C^i\) be the cohomology of the
cone of restriction from global cochains to local cochains
of \(Q_v\), shifted by \(-1\). In degree one this keeps a
global cocycle and local trivializations modulo \(b_v\).
If \(c_w\) denotes conjugation at a real place, the desired
bound is
\begin{equation}\label{eq:fitted-relative-euler}
 \dim H_C^2+\dim H_C^0+\sum_{v\mid\ell}d_v\dim b_v
 \ \le\ \dim H_C^1+\sum_{w\mid\infty}\dim W^{c_w}.
\end{equation}

\end{proposition}

\begin{proof}
\proofstep{Comparison with the arithmetic cohomology groups.}
To compare with global duality, first enlarge the set of
allowed ramification of \(L^+\) from the \(\ell\)-places to a
finite set \(S'\) including the ramification of \(\eta\) and
\(P\). Denote the unrestricted group on \(S'\) by \(\widetilde G\).
Its map to \(G\) has kernel \(J\) normally generated by the
images of \(J_q=I_q\cap\ker\eta\cap\ker P\) at
\(q\in S'\), \(q\nmid\ell\). The quotient \(I_q/J_q\) has order
prime to \(\ell\). Thus a class of \(H^2(G,W)\) has inflated
restriction zero at each such \(q\). Its local inflation
factors via the decomposition group modulo \(J_q\), an extension
of the procyclic unramified quotient by a finite group of
order prime to \(\ell\), whose \(\ell\)-cohomological dimension
is at most one.

Restriction gives an injection
\[
 \operatorname{Hom}_{\rm cont}(J,W)^{\widetilde G}
 \longrightarrow \bigoplus_{q\nmid\ell}\operatorname{Hom}_{D_q}(J_q,W).
\]
The restrictions determine a homomorphism on the closed
normal generators. The dimension of the local summand is
\(h^2(D_q,W)\), with \(h^i\) denoting dimension of continuous
cohomology. In fact wild inertia has order prime to \(\ell\)
in its finite quotients, and the pro-\(\ell\) part of tame
inertia is \(\mathbb Z_\ell(1)\), unchanged by passing to
the prime-to-\(\ell\) index subgroup. Equivariance asks for the
Frobenius-fixed space of \(W(-1)^{I_q}\). Local duality counts
the coinvariants of the same space under Frobenius, of the
same dimension. Averaging over \(I_q/J_q\) identifies inertia
invariants and coinvariants in this comparison.

\proofstep{The relative tangent and obstruction terms.}
Write \(a=h^1(\widetilde G,W)-h^1(G,W)\), and let \(T\)
denote the invariant homomorphism space on \(J\) just
considered. The relevant five-term sequence is
\[
 0\to H^1(G,W)\to H^1(\widetilde G,W)\to T
 \to H^2(G,W)\to H^2(\widetilde G,W).
\]
The image of \(H^1(\widetilde G,W)\) in \(T\) has dimension \(a\);
therefore the kernel at \(H^2(G,W)\) has dimension \(\dim T-a\).
Let \(K_2\) be the kernel of the map from \(H^2(G,W)\) to
\(\bigoplus_{v\mid\ell}H^2(D_v,W)\).
The image of \(K_2\) is locally zero at every place of
\(S'\), including the deleted primes by the inflation
argument.

In particular \(H^2(G,W)\) is finite-dimensional,
and the Poitou--Tate sequence on \(S'\) gives
\[
 \dim K_2 \le \dim T-a+h^2(\widetilde G,W)
                         -\sum_{v\in S'}h^2(D_v,W).
\]
The global dual invariant term vanishes by hypothesis.
Use the bound on \(\dim T\) and the global Euler characteristic
for \(\widetilde G\). It follows that
\[
 \dim K_2+\sum_{v\mid\ell}h^2(D_v,W)
        +h^0(G,W)+d\dim W
 \le h^1(G,W)+\sum_{w\mid\infty}\dim W^{c_w}.
\]
Passing from \(W\) to \(Q_v\) in each local target makes the
coefficient map on \(H^2\) surjective, since local
\(\ell\)-cohomological dimension is two. The kernel dimension
increases by at most the summed drop in the local dimensions.
The last inequality therefore also holds with the new kernel
dimension \(\kappa_2\) and \(h^2(D_v,Q_v)\).

The cone sequence gives
\begin{align*}
 \dim H_C^2-\dim H_C^1+\dim H_C^0
   &=\kappa_2-h^1(G,W)+h^0(G,W)\\
   &\quad+\sum_{v\mid\ell}\bigl(h^1(D_v,Q_v)-h^0(D_v,Q_v)\bigr).
\end{align*}
By the local Euler formula the summand is
\(h^2(D_v,Q_v)+d_v\dim Q_v\). Use the preceding global
inequality and \(\sum_{v\mid\ell}d_v=d\) to obtain
\eqref{eq:fitted-relative-euler}. This gives the relative
relation count on the same local groups and weights that
give the marked first cohomology bounds.
\cite{milne-duality}\label{cor:ordinary-fitted-arithmetic-4}

\end{proof}

\paragraph{The three auxiliary level actions.}
It remains to realize the auxiliary level actions with the prescribed local
characters. The global separating element supplies, by Chebotarev on
the finite quotient recording \(L,P,\bar\epsilon\), degree-one
split places \(v_0\) of \(L^+\) with \(P=1\) on the whole
decomposition group and with nontrivial mod-\(\ell\) cyclotomic
value at Frobenius. Indeed the characters are unramified at
these primes, and the chosen Frobenius fixes \(L\) itself.
Avoid the primes of \(2M\operatorname{disc}(L)\) and any
previous choice. A fixed finite scalar twist keeps all
residual readings on the common scalar line there. If the common residual scalar is \(a\ne0\), the scalar after the cyclotomic twist is different at the chosen Frobenius, because
\[
 \bar\epsilon(\mathrm{Frob}_{v_0})\ne1\quad\Longrightarrow\quad a\bar\epsilon(\mathrm{Frob}_{v_0})\ne a.
\] The same choice makes the local compact-group
condition at \(v_0\) neat, as in
Paragraph~\ref{par:ordinary-functions}.

Use the appropriate scalar twist for each of the rank-two
and two rank-three levels in Section~\ref{sec:linked-seeds}.
At a place above the tame cutoff the inertia type consists
of the diagonal Teichm\"uller characters after this scalar
twist. The same type and the same ordered \(\ell\)-adic
characters are retained at every auxiliary depth of that
level. The two rank-three levels have different
distinguished places. At each of them, the other
distinguished place remains outside the ramification set
of that level. The two actions are compared separately at
their own hyperspecial starting modules before contraction
to the common frame.

\begin{proposition}[Auxiliary actions at the three actual levels]
\label{prop:fitted-three-level-actions}
Under the fitted level conditions, let \(M^{(3)},M^{(0)},M^{(2)}\) denote the modules for the
rank-three primary, rank-three auxiliary and rank-two actions, respectively.
The annihilators for these level actions, computed at empty auxiliary set, are
\[
  \operatorname{Ann}(M^{(3)})=J_3,\qquad
  \operatorname{Ann}(M^{(0)})=J_0,\qquad
  \operatorname{Ann}(M^{(2)})=J_2 .
\] These actions provide the fixed spaces
in the patching argument. The descended containment on the original
host acquires no factor from the fitted degree or number of
Taylor--Wiles places.
\end{proposition}

\begin{proof}
\proofstep{The finite groups at the actual Hecke level.}
Here is the identification at a single level. Enlarge the
allowed ramification set by the selected Taylor--Wiles
places. If \(S\) is the original set of allowed places and \(Q\) the auxiliary set, the quotient at the level in question is
\[
 G_{S\cup Q}\big/\overline{\langle\!\langle I_v:v\in Q\rangle\!\rangle}
      \simeq G_S .
\]
Here the denominator is the closed normal subgroup generated by the images of the added inertia groups.

At hyperspecial compact level at all
new places, the function space and ordinary operator
are exactly the original ones. The ordinary residual blocks
also agree. Indeed the existing operators and those retained
at common spherical places agree, and a classical point
belongs to the residual block exactly when its residual
Frobenius polynomials equal the fixed diagonal polynomials.
Frobenius elements at the common spherical places are dense
up to conjugacy in the quotient. The continuous residual determinant is determined on the quotient by those conjugacy classes. Indeed two such determinants with identical characteristic polynomials on the dense Frobenius set have
\[
 D_1(1-Xg)=D_2(1-Xg)\qquad(g\text{ in the quotient})
\]
by continuity of each coefficient. The local weight character is fixed. Separating classical weights
and finite freeness give the same equality on the integral
ordinary modules. This identifies the empty-auxiliary
annihilator with the Hecke ideal of that level.

At the varying auxiliary levels take the directed union
of the modules of Paragraph~\ref{par:ordinary-functions}.
Right translation on this union is the smooth action
of the product of the auxiliary general linear groups.
The universal determinant and weight ring acts compatibly
on the fixed spaces. For auxiliary levels \(K'\subset K\), smaller level subgroup means a larger fixed space. On the underlying module \(M\) the transition is the inclusion
\[
 i_{K,K'}:M^K\hookrightarrow M^{K'},\qquad v\longmapsto v .
\] For an element \(d\) of the determinant and weight ring form
\[
 A_d=i_{K,K'}\,d_K-d_{K'}\,i_{K,K'} .
\]
At every separating classical specialization, the Frobenius polynomials and the ordered weight characters identify the two determinant actions, so \(A_d\) specializes to zero. On every separating classical specialization \(x\) the defect is zero.

Since these specializations separate the projected ordinary module, an element of the image of this defect has all its specializations zero, so
\[
 \operatorname{im} A_d\subseteq\bigcap_x\ker(\text{specialization at }x)=0 .
\] The same comparison after a change of level by conjugation intertwines the determinant action with translation. Thus these local actions and the
determinant action commute on the directed union.
At the tame level the same modules are projective over
the group ring of the finite \(\ell\)-group quotient by
the argument for the neat level. The augmentation trace
and the two maps with product of split-root factors in
Paragraph~\ref{par:ordinary-level} apply to these fixed
spaces. The integer controlling the local nilpotent relations
depends only on \(n\), and no auxiliary depth or factor is
charged to the final host trace product.

\proofstep{Keeping the original coefficient field in the comparison.}
In the rank-three linked comparison, the primary level
uses the projected rank-three multiplicity computation.
The auxiliary rank-three level and the rank-two level use
the same ordinary construction without the separate
rank-three multiplicity computation. The resulting annihilator ideals, at the three corresponding levels in order, are
\[
 J_3\quad\text{(primary rank three)},\qquad J_0\quad\text{(auxiliary rank three)},\qquad J_2\quad\text{(rank two)} .
\] The patching argument of Section~\ref{sec:marked-patching} on these auxiliary fixed spaces gives ideal containments on the indicated curves. The valuation bound after descent to the
rational host is independent of the dimension bounds on the
fitted field and the number of Taylor--Wiles places.
\label{cor:ordinary-fitted-arithmetic-5}
\end{proof}
